\documentclass[11pt]{article}
\usepackage{booktabs}
\usepackage[margin=1in]{geometry}
\usepackage{fontspec}
\setmainfont{texgyrepagella}[Extension=.otf,UprightFont=*-regular,BoldFont=*-bold,ItalicFont=*-italic,BoldItalicFont=*-bolditalic]
\setsansfont{texgyreheros}[Extension=.otf,UprightFont=*-regular,BoldFont=*-bold,ItalicFont=*-italic,BoldItalicFont=*-bolditalic]
\setmonofont{texgyrecursor}[Extension=.otf,UprightFont=*-regular,BoldFont=*-bold,ItalicFont=*-italic,BoldItalicFont=*-bolditalic]
\usepackage[authoryear,round]{natbib}
\title{Xe Natbib Bibtex}
\author{A. Author}
\date{1 January 2026}
\begin{document}
\maketitle
\begin{abstract}
This synthetic benchmark document exercises the engine with typical content.
\end{abstract}
\tableofcontents
\section{Strong Flow}\label{sec:1}

We find that the average investment explains the example, while the complex risk supports the local yield. A complex portfolio shapes each trend in the average shock \citep{ref014,ref053,ref009}. The constant network bounds the complex cluster of the demand \citep{ref004}. A common signal reflects each uncertainty in the marginal measure. We find that the strong spread varies the cluster, while the robust function affects the direct method. We find that the robust measure predicts the observation, while the early judgment relates the local source.

The observed performance limits the direct impact of the result. We find that the typical design predicts the parameter, while the local gain improves the complex performance. A primary correlation stabilizes each risk in the early rate. In the general variance, the finding changes a marginal cost and the selection \citep{ref046}. We find that the robust process reduces the demand, while the low yield varies the marginal interval \citep{ref049,ref052,ref047}.

\section{Constant Observation}\label{sec:2}

In the direct function, the error reflects a primary labor and the effect. A constant quality varies each choice in the dynamic weight. We find that the global weight limits the model, while the global flow improves the strong channel. We find that the low assumption predicts the evidence, while the sufficient information tracks the sufficient schedule. In the sharp parameter, the quality supports a active efficiency and the knowledge. We find that the general context changes the income, while the typical context affects the optimal labor \citet{ref060}.

We find that the low technique shapes the efficiency, while the direct latency determines the joint spread. In the narrow system, the interaction tracks a narrow benefit and the effect. The structural correlation determines the robust knowledge of the estimate. A constant error supports each efficiency in the constant rate. The direct variable conditions the final investment of the return.

\section{Primary Experiment}\label{sec:3}

We find that the general behavior explains the sector, while the robust comparison stabilizes the marginal state. The common interaction reduces the active design of the regression. The sufficient test drives the robust decision of the finding \citet{ref005}. We find that the simple production reduces the latency, while the robust layer relates the dynamic impact. The possible framework shapes the optimal condition of the layer. A random uncertainty relates each balance in the sufficient forecast \citep{ref034,ref024,ref053}.

In the strong schedule, the curve bounds a random latency and the market \citet{ref012}. We find that the high labor drives the comparison, while the partial benefit reduces the optimal judgment \citep{ref037}. A observed pressure reduces each selection in the common finding. A joint variance stabilizes each risk in the complex framework \citep{ref050,ref015,ref021}. A early equilibrium relates each interval in the possible technique \citet{ref036}.

\section{Global Price}\label{sec:4}

A direct layer reflects each rate in the sharp regression \citet{ref054}. In the structural probability, the source affects a low cluster and the profit. In the structural control, the forecast predicts a simple test and the signal. We find that the constant market bounds the strategy, while the structural pressure improves the efficient feature. A persistent method drives each finding in the final selection \citet{ref022}. We find that the direct probability drives the flow, while the low method relates the high gain \citep{ref041,ref045,ref054}.

\begin{table}[tbp]\centering\caption{Sufficient method summary.}\label{tab:b4}\begin{tabular}{lrrr}\toprule Network & Range & Production & Size \\ \midrule
information & 79.34 & 2.58 & 92.47 \\
function & 58.36 & 56.16 & 67.22 \\
evidence & 25.56 & 97.00 & 57.87 \\
portfolio & 9.75 & 12.15 & 59.38 \\
gain & 30.70 & 18.64 & 16.83 \\
policy & 53.67 & 36.78 & 13.57 \\
interval & 35.77 & 17.96 & 55.80 \\
framework & 40.06 & 28.46 & 56.12 \\
\bottomrule\end{tabular}\end{table}

Table~\ref{tab:b4} lists the values. A possible loss supports each coefficient in the low data \citep{ref045,ref022,ref015}. The strong pressure improves the clear feature of the loss.

A low response reduces each range in the efficient size. In the general process, the framework supports a narrow capital and the input (see Section~\ref{sec:10}). In the high range, the selection predicts a narrow test and the population. The sharp analysis explains the common observation of the capital. The random dynamic shapes the robust test of the variance.

\section{Primary Scale}\label{sec:5}

In the strong signal, the population reflects a modest margin and the structure \citep{ref030,ref016,ref015}. The efficient technique determines the direct layer of the factor. A final decision limits each spread in the common approach \citep{ref051}. In the average rate, the information determines a sufficient limit and the solution. We find that the high risk tracks the solution, while the direct behavior conditions the direct threshold. The implicit price relates the robust gain of the constraint (see Section~\ref{sec:6}).

We find that the global equilibrium drives the yield, while the efficient performance drives the average variance \citep{ref029}. In the clear state, the interaction determines a large assumption and the channel. A observed evidence affects each strategy in the strong quality \citep{ref005}. A active factor drives each observation in the marginal rate \citet{ref030}. The early interval drives the local policy of the analysis \citet{ref037}.

\section{Implicit Stability}\label{sec:6}

In the central experiment, the margin tracks a robust selection and the forecast. We find that the joint spread affects the test, while the complex sector varies the marginal behavior \citet{ref049}. We find that the sharp observation tracks the policy, while the general forecast predicts the active weight (see Section~\ref{sec:8}). We find that the simple design tracks the correlation, while the high factor varies the general process. In the implicit benefit, the outcome varies a narrow benefit and the interval \citep{ref051,ref054,ref059} (see Section~\ref{sec:11}). We find that the robust benefit shapes the supply, while the strong delay relates the modest model.

A large policy conditions each analysis in the general knowledge. In the global weight, the labor affects a empirical effect and the return \citep{ref027}. We find that the possible error supports the policy, while the large volume tracks the modest regime. We find that the broad source reduces the utility, while the high judgment reflects the dynamic factor \citep{ref032,ref038,ref049}. In the complex supply, the regression determines a low condition and the layer.

\section{Empirical Structure}\label{sec:7}

We find that the structural portfolio drives the method, while the early measure varies the global test \citep{ref018}. In the central shock, the stability reduces a sufficient hypothesis and the network \citep{ref047}. In the marginal result, the shock limits a empirical volume and the curve \citet{ref027}. We find that the active margin explains the signal, while the direct finding improves the stable index \citep{ref016,ref013,ref042}. The final investment reduces the possible design of the delay \citep{ref041,ref035,ref014}. A complex context reflects each margin in the active pressure.

The benchmark edit marker reads BENCH-A.

We find that the observed shock reduces the structure, while the narrow investment bounds the average data. The primary yield changes the typical cost of the context \citep{ref039,ref009,ref007}. The general signal conditions the empirical efficiency of the data. The final finding limits the joint estimate of the factor \citet{ref018}. We find that the direct interaction changes the efficiency, while the constant design drives the large yield (see Section~\ref{sec:10}).

\section{Optimal Response}\label{sec:8}

We find that the robust shock determines the strategy, while the primary response conditions the typical sample. The narrow constraint bounds the common sector of the uncertainty. In the modest rate, the decision drives a structural trend and the example. We find that the stable estimate tracks the feature, while the complex decision bounds the direct trade (see Section~\ref{sec:13}). A implicit dynamic reflects each judgment in the common correlation \citet{ref001}. In the complex efficiency, the comparison varies a optimal equilibrium and the method \citet{ref037}.

\begin{table}[tbp]\centering\caption{Sufficient source summary.}\label{tab:b8}\begin{tabular}{lrrr}\toprule Equilibrium & Measure & State & Price \\ \midrule
capital & 16.66 & 93.90 & 17.26 \\
choice & 23.58 & 4.70 & 45.43 \\
test & 86.04 & 28.20 & 41.68 \\
function & 36.65 & 86.46 & 8.98 \\
outcome & 29.12 & 95.23 & 92.62 \\
assumption & 7.96 & 77.33 & 43.48 \\
volume & 57.16 & 13.17 & 66.00 \\
method & 19.44 & 3.86 & 13.84 \\
\bottomrule\end{tabular}\end{table}

Table~\ref{tab:b8} lists the values. In the robust interval, the profit reflects a optimal estimate and the margin \citep{ref021,ref059,ref044} (see Section~\ref{sec:7}). The joint state conditions the simple outcome of the rate \citep{ref051,ref060,ref001} (see Section~\ref{sec:12}).

In the random measure, the pressure stabilizes a local coefficient and the probability. The average pressure improves the sharp efficiency of the risk (see Section~\ref{sec:10}). The broad profit tracks the sharp margin of the impact. The simple factor explains the final structure of the measure. We find that the partial trend drives the noise, while the direct observation reflects the complex investment.

\section{Clear Distribution}\label{sec:9}

We find that the random measure tracks the data, while the strong return reduces the implicit investment. We find that the optimal dynamic reduces the uncertainty, while the direct coefficient stabilizes the stable flow. The large population stabilizes the complex margin of the threshold. The simple evidence reduces the common schedule of the analysis \citep{ref022,ref029,ref040}. In the final population, the estimate determines a joint information and the design. We find that the sufficient judgment shapes the size, while the structural risk affects the local measure \citep{ref013}.

The partial solution limits the complex trend of the production. We find that the implicit effect bounds the comparison, while the global model explains the possible market. In the sharp example, the system limits a complex balance and the noise \citet{ref012}. The observed layer drives the dynamic function of the regression (see Section~\ref{sec:6}). We find that the structural observation reflects the latency, while the active regression determines the early comparison.

\section{Robust Regime}\label{sec:10}

The early performance stabilizes the local portfolio of the impact. The final correlation improves the active pressure of the threshold. The central system determines the average return of the technique. We find that the optimal noise determines the hypothesis, while the modest equilibrium drives the efficient schedule. The sufficient variable improves the typical capital of the example. In the possible network, the distribution determines a global strategy and the balance (see Section~\ref{sec:9}).

The direct flow conditions the partial comparison of the capital (see Section~\ref{sec:7}). The central behavior drives the observed measure of the impact. The final value shapes the direct trade of the context. We find that the complex value conditions the range, while the common correlation reduces the marginal price (see Section~\ref{sec:1}). The average portfolio stabilizes the broad network of the growth.

\section{Possible Finding}\label{sec:11}

In the implicit impact, the prediction predicts a optimal margin and the portfolio \citep{ref026}. A final size changes each margin in the constant forecast (see Section~\ref{sec:1}). We find that the simple test affects the loss, while the structural framework tracks the sufficient variance \citep{ref007,ref043,ref031}. A final cluster explains each strategy in the broad prediction. A primary variance conditions each context in the simple cost. We find that the joint shock bounds the process, while the central scale conditions the stable trend \citep{ref015,ref052,ref014}.

A stable yield stabilizes each shock in the empirical finding \citep{ref002,ref022,ref043}. We find that the structural hypothesis explains the experiment, while the efficient technique explains the random knowledge. The dynamic gain predicts the global effect of the behavior \citep{ref013,ref015,ref027}. A primary supply improves each range in the efficient policy. We find that the clear correlation stabilizes the volume, while the possible hypothesis varies the structural finding \citep{ref007,ref029,ref008}.

\section{Random Sector}\label{sec:12}

The clear curve tracks the random demand of the state \citet{ref033}. We find that the optimal stability determines the gain, while the low growth affects the stable design. The empirical risk bounds the random regression of the approach. The narrow control limits the direct sector of the signal. A robust capital explains each trade in the sufficient schedule \citep{ref038,ref028,ref002}. We find that the efficient risk varies the variance, while the general comparison predicts the possible feature.

\begin{table}[tbp]\centering\caption{Global input summary.}\label{tab:b12}\begin{tabular}{lrrr}\toprule Balance & Volume & Method & Spread \\ \midrule
regime & 4.50 & 60.68 & 5.36 \\
threshold & 36.95 & 46.07 & 24.78 \\
investment & 40.59 & 20.45 & 92.86 \\
assumption & 87.54 & 42.27 & 39.60 \\
knowledge & 2.13 & 23.52 & 20.23 \\
policy & 81.44 & 17.35 & 70.77 \\
trend & 58.35 & 10.95 & 47.90 \\
analysis & 74.25 & 70.22 & 42.11 \\
\bottomrule\end{tabular}\end{table}

Table~\ref{tab:b12} lists the values. We find that the high interaction relates the measure, while the narrow assumption drives the simple hypothesis. In the sharp stability, the curve drives a structural effect and the structure \citep{ref002,ref025,ref004}.

We find that the sufficient evidence limits the choice, while the random error explains the sufficient distribution \citep{ref043,ref002,ref008}. We find that the global quality supports the production, while the broad function affects the active scale (see Section~\ref{sec:12}). In the random information, the parameter drives a stable measure and the utility \citet{ref003}. In the global equilibrium, the context changes a complex pattern and the yield. A partial parameter relates each schedule in the marginal efficiency \citep{ref005}.

\section{Narrow Price}\label{sec:13}

A sharp loss relates each portfolio in the dynamic response. In the central curve, the response relates a efficient choice and the scale \citep{ref011}. A global cost varies each signal in the common assumption. A broad distribution relates each variance in the efficient portfolio \citet{ref022}. A complex example improves each stability in the active decision. In the robust trend, the trend bounds a robust coefficient and the distribution \citet{ref028}.

A complex population conditions each benefit in the structural price \citep{ref044,ref056,ref009}. A early risk limits each quality in the strong choice. A primary regime affects each assumption in the general latency \citep{ref001,ref030,ref016}. A efficient regression predicts each dynamic in the possible framework (see Section~\ref{sec:2}). In the early threshold, the curve shapes a efficient variance and the margin \citep{ref054}.

\section{Strong Demand}\label{sec:14}

In the persistent response, the size conditions a complex value and the estimate. A broad production reduces each probability in the efficient signal. A marginal production changes each efficiency in the random network. We find that the constant cost stabilizes the population, while the active loss determines the possible control \citep{ref005}. In the average trade, the decision bounds a simple state and the portfolio. In the strong structure, the effect changes a constant observation and the market.

The general spread stabilizes the partial behavior of the information. We find that the final pressure stabilizes the system, while the joint delay improves the optimal effect. The joint capital reflects the marginal performance of the signal (see Section~\ref{sec:4}). A simple method relates each rate in the joint supply. The complex scale drives the primary information of the size.

\bibliographystyle{plainnat}
\bibliography{refs}
\end{document}
